\begin{table}[!tbh]
\centering\scriptsize\setlength{\tabcolsep}{3pt}
\begin{tabular}{|l|p{.40\linewidth}|p{.19\linewidth}|l|}
\hline
& \textbf{Intervention} & \textbf{Evaluation} & \textbf{Status} \\
\hline
E1 & search method at matched candidate diversity & 300 tasks, 3 seeds & confirmatory \\ \hline
E2 & 400 steps of GRPO with an exploring sampler & 300 tasks, 3 seeds & confirmatory \\ \hline
E3 & supervision on correct depth-2 and depth-3 programs & $4 \times 300$, 3 seeds & confirmatory \\ \hline
E4 & removal of three operation pairs from supervision & $4 \times 250$, 3 seeds & confirmatory \\ \hline
E5 & equal-budget control, motif-present control, APPLY mixture, depth 4, verifier noise & reused sets, 125 new tasks & post-hoc \\ \hline
\end{tabular}
\caption{The five studies and their status, stated before any result is shown. E1
to E4 were preregistered, with the primary endpoint, the decision rule and the
frozen set fixed in advance. E4 was added by an amendment whose position in time
relative to the opening of E3's data cannot be established from the artifacts, so it
is a disclosed extension and not an independently preregistered test. E5 collects
boundary conditions on sets
that E3 and E4 had already opened, together with one newly generated depth-4
split, and is therefore reported as post-hoc. The verifier levels of E5 use one
training seed each and the rest of E5 uses three. The claims of E1 to E4 stand
without E5.}
\label{tab:design}
\end{table}
